\chapter*{Kurzfassung}
Diese Bachelorarbeit untersucht strategische Architekturentscheidungen im Bereich moderner IT-Infrastrukturen. 
Im Mittelpunkt steht die Frage, in Abhängigkeit welcher strategischen, wirtschaftlichen und organisatorischen Faktoren 
Unternehmen zwischen On-Premise-, Cloud- oder hybriden Architekturkonzepten wählen sollten. 

Ziel der Arbeit ist die Entwicklung eines Bewertungsmodells, das betriebliche Entscheidungsträger bei der Auswahl 
geeigneter Infrastrukturansätze unterstützt. Grundlage bildet eine systematische Analyse wissenschaftlicher Literatur, 
anerkannter Frameworks und aktueller Praxisbeispiele aus dem Bereich IT-Management und Cloud-Governance. 
Das daraus abgeleitete Modell soll als konzeptionelle Entscheidungshilfe dienen und die Integration 
strategischer, ökonomischer und organisatorischer Kriterien in Architekturentscheidungen ermöglichen. 

\vfill
\textbf{Schlüsselwörter:} Strategische IT-Architekturen, Cloud Computing, On-Premise, Hybrid, Entscheidungsmodell, Wirtschaftsinformatik

